% ------------------------------------------------------------------
% Log of the experiment "learned reconstruction of the ancestor of
% two tandem repeat sequences" (17 November 2026).
% Notation follows Matroud, Tuffley and Hendy (TCBB, 2021) and the
% MPE manuscript: A (ancestor), D_1, D_2 (descendants), \hat{A},
% n(A), e, \varepsilon, d_est, d_hist, \delta_rel, the (s,t) grid.
% ------------------------------------------------------------------
\documentclass[11pt,a4paper]{article}
\usepackage[margin=2.5cm]{geometry}
\usepackage{amsmath,amssymb}
\usepackage{booktabs}
\usepackage{xcolor}
\usepackage{hyperref}
\newcommand{\todo}[1]{\textcolor{red}{[\textbf{To confirm:} #1]}}
\newcommand{\dest}{d_{\mathrm{est}}}
\newcommand{\dhist}{d_{\mathrm{hist}}}
\newcommand{\drel}{\delta_{\mathrm{rel}}}
\newcommand{\Ahat}{\hat{A}}

\title{Learned reconstruction of the ancestor of two tandem repeat sequences:\\ log of the experiment}
\author{Atheer Matroud}
\date{17 November 2026}

\begin{document}
\maketitle

\section{Aim}

In the manuscript \emph{Tandem repeats ancestor reconstruction and distance} the ancestor $A$ of two tandem repeat sequences $D_1$ and $D_2$ is reconstructed by an alignment-based procedure. In this experiment a model is instead trained on simulated families $(A;D_1,D_2)$, in which the ancestor is known, and is then asked to reconstruct the ancestor of families that it has not seen. The reconstructed ancestor $\Ahat$ is compared with the original ancestor $A$ by the measures of the manuscript, and the result is compared with that of the alignment-based procedure.

\section{Decisions taken}

The following were confirmed on 17 November 2026.
\begin{enumerate}
\item The families cover the whole range of the manuscript, $2\le s\le 20$ and $1\le t\le 20$.
\item Three candidate models are compared and the best is kept.
\item The data set consists of 10\,000 families, each with two descendants.
\item Training is carried out on a Colab GPU, and the work is arranged to run in parallel.
\item In addition to the accuracy of $\Ahat$, the distance estimate $\dest(\Ahat,D_i)$ is compared with $\dhist(A,D_i)$.
\item The comparison is made both with the original ancestor $A$ and with the result of the manuscript.
\item The code is kept in the folder \texttt{AI} of the repository \url{https://github.com/AtheerMatroud/TandemRepeats}.
\end{enumerate}

\section{Data}

\subsection{Generation}

Every family is generated by the generator of the manuscript (\texttt{tr\_aligner.simulate.TCBBSimulator}), without change: an initial sequence of $50s$ copies of a motif of length 11 chosen uniformly at random; an ancestral sequence $A$ formed by tandem repeat edit events until it carries 15 motif variants; and two descendants $D_1$ and $D_2$, each formed independently by a further $5t$ events with the rate ratio \emph{deletion} : \emph{substitution} : \emph{duplication} of $1:5:20$, duplicated segments of at most 5 copies and deleted segments of at most 3 copies. For each of the 10\,000 families the pair $(s,t)$ is drawn uniformly at random from the $19\times 20$ grid, so that every part of the grid is represented by about 26 families. The weight $\dhist(A,D_i)$ of the true history is recorded for each descendant.

\subsection{Training, validation and test sets}

The families are divided into a training set (80\%), a validation set (10\%) and a test set (10\%). The division is made separately for the families of low event density, $5t<\sqrt{2n(A)\ln 2}$, and those of high event density, so that the three sets have the same proportion of each.

\subsection{The families of the manuscript}

The 380 families with $k=2$ of the manuscript are regenerated from the same seed as a second test set. It was verified that the regenerated families are those of the manuscript: $n(A)$ and $\dhist(A,D_1)$ agree with the recorded result files in all 380 cells, and the alignment-based procedure applied to the regenerated families of the cells $(s,t)=(5,3)$ and $(12,10)$ returned the recorded $n(\Ahat)$, $e(A,\Ahat)$ and $\dest$. The recorded results of the manuscript for these families are kept in \texttt{data/mpe\_k2\_results.csv}; the summary computed from that file reproduces the figures of the manuscript for $k=2$ (mean $\varepsilon$ of 2.7\% for the 92 low density families, 20.2\% for the 288 high density families and 16.0\% overall).

\section{Representation of a family}

A motif variant has no meaning outside its own family, since the motif is chosen at random in every family. The variants of a family are therefore numbered in the order in which they first appear in $D_1$ and then in $D_2$, and every symbol presented to a model is accompanied by the nucleotides of its variant, so that the model can see which variants differ by a single substitution. A variant of $A$ that occurs in neither descendant cannot be produced by any of the models.

The training families are augmented by exchanging $D_1$ and $D_2$ and, for the two models that produce $A$ symbol by symbol, by reversing the three sequences.

\section{Candidate models}

\begin{description}
\item[Transformer.] An encoder-decoder Transformer (4 encoder and 4 decoder layers, width 256). The input is $\tau_v(D_1)$, a separator and $\tau_v(D_2)$; the output is $\tau_v(A)$, produced one symbol at a time.
\item[BiLSTM.] A bidirectional LSTM encoder and an LSTM decoder with attention, with the same input and output.
\item[Alignment-guided model.] $\tau_v(D_1)$ and $\tau_v(D_2)$ are first aligned by a symmetric global alignment (match 0, substitution 20 per differing nucleotide, gap 10 per motif copy). A bidirectional LSTM (3 layers, width 256) then reads the columns of the alignment and decides, for every column, whether $A$ has the symbol of $D_1$, the symbol of $D_2$, or no copy at that position. $\Ahat$ is spelled by the columns that are kept. For training, the decision for each column is obtained from the true ancestor by a longest common subsequence calculation between $\tau_v(A)$ and the columns.
\end{description}

\todo{The alignment used inside the third model is a plain symmetric alignment and not the asymmetric alignment of the TCBB paper. It was chosen because it takes milliseconds per family, whereas the asymmetric alignment would need many hours for 10\,000 families. The weights are those of the sibling scoring of the manuscript without the duplication term. The names ``alignment-guided model'' and ``column'' are not terms of the TCBB paper and are to be confirmed or replaced before any of this enters a manuscript.}

The model is selected on the validation set only, by the mean relative error $\varepsilon(A,\Ahat)=e(A,\Ahat)/n(A)$.

\section{Evaluation}

\subsection{Accuracy of the reconstructed ancestor}

As in the manuscript, $\tau_v(\Ahat)$ is compared with $\tau_v(A)$ by the symbol edit distance $e(A,\Ahat)$. The number of families recovered exactly ($e=0$), the mean and median of $e$, the mean of $\varepsilon(A,\Ahat)$ and the mean of $n(\Ahat)-n(A)$ are reported for the low density families, the high density families and all families.

\subsection{Comparison with the result of the manuscript}

Two comparisons are made.
\begin{enumerate}
\item On the 1000 families of the test set, the alignment-based procedure of the manuscript is run on the same families, and the two reconstructions are compared family by family (mean difference of $\varepsilon$, proportion of families in which each is better, Wilcoxon signed-rank test).
\item On the 380 families with $k=2$ of the manuscript, the model is compared with the results recorded for the manuscript, cell by cell.
\end{enumerate}

\subsection{Distance estimates}

For each descendant, $\dest(\Ahat,D_i)$ is calculated by the asymmetric alignment from the ancestor reconstructed by the model, under the match-reward scoring and the edit-weight scoring, and
\[
\drel=\frac{\dest-\dhist(A,D_i)}{\dhist(A,D_i)}
\]
is tabulated (sign counts and summary statistics, in the format of the manuscript) beside the same quantity for the ancestor reconstructed by the alignment-based procedure and for the true ancestor.

\section{Notebooks}

\begin{center}
\begin{tabular}{lll}
\toprule
notebook & content & hardware \\
\midrule
\texttt{01\_generate\_dataset} & the 10\,000 families, the three sets, the 380 families of the manuscript & CPU \\
\texttt{02\_baseline\_alignment} & alignment-based procedure on the test set & CPU \\
\texttt{03\_train\_models} & training of the three candidate models & GPU \\
\texttt{04\_evaluate} & test, comparisons, distance estimates, tables and figures & GPU, then CPU \\
\bottomrule
\end{tabular}
\end{center}

All notebooks write to the Drive folder \texttt{tandem\_repeats/ai\_experiment}. Work in parallel is arranged in three ways.
\begin{enumerate}
\item Every CPU stage uses all the cores of its session and proceeds in small independent parts; a part that is finished is never recomputed, so that an interrupted session is simply run again.
\item A notebook can be opened in several sessions at once, each with its own value of \texttt{PART}; the sessions share the Drive folder.
\item The three models can be trained in three GPU sessions at the same time, one model in each, and notebook 2 needs no GPU and can run while the models train.
\end{enumerate}

\section{Checks carried out before delivery}

\begin{itemize}
\item The four notebooks were executed from start to finish on a reduced data set (200 families) on a CPU, including the training of two models, the selection, the test, the comparison tables and the distance estimates.
\item The summary of the recorded results of the manuscript reproduces its figures for $k=2$ (Section 3.3).
\item The columns of the symmetric alignment can express the true ancestor almost completely: in 30 families taken across the grid, the best possible choice of columns gave $e(A,\Ahat)=0$ in 28 families and $e=2$ in the other two. The alignment itself therefore does not limit the accuracy of the third model.
\item A small version of the third model (0.4 million parameters) was trained on a CPU on 600 families with $s\le 6$ and tested on 60 further families of the same kind. The mean of $\varepsilon(A,\Ahat)$ was 15.5\%, against 31.1\% for the alignment-based procedure on the same 60 families; the mean of $n(\Ahat)-n(A)$ was 7.0 copies against 54.6. This is a preliminary check on a small sample and on short ancestors only; it is not a result of the experiment.
\end{itemize}

\section{Points to bear in mind when the results are read}

\begin{itemize}
\item A trained model learns the rates of the generator, in particular that a copy present in one descendant only is far more often a duplication in that descendant than a deletion in the other. The alignment-based procedure makes no such assumption. An improvement on simulated data is therefore expected to transfer to real data only in so far as the real process resembles the simulated one, and this should be stated with any result.
\item The ancestors of the generator consist largely of copies of a single variant, the other 14 variants being rare. Where both descendants have gained copies within a run of the same variant, the number of copies in $A$ cannot be determined exactly from $D_1$ and $D_2$, and a model can do no better than estimate it.
\item The two models that produce $A$ symbol by symbol must learn to copy sequences of up to about 1000 symbols from 8000 examples, and may not reach the accuracy of the third model. They are included because the comparison of the three was requested.
\end{itemize}

\section{Results}

\todo{To be completed from the files written by notebook 4 (\texttt{table\_accuracy\_test.csv}, \texttt{table\_accuracy\_mpe.csv}, \texttt{table\_delta\_rel.csv} and the figures) once the notebooks have been run.}

\end{document}
